Der Körper bewege sich ohne Luftwiderstand auf einem Kreis mit dem Erdradius $R = \rErO$ (Äquatorradius, FoTaBe S.~212). Die Gewichtskraft liefert die Zentripetalkraft:
\begin{align}
 m\,g &= \frac{m\,v^2}{R}
\end{align}
Die Masse kürzt sich, und es folgt
\begin{align}
 v &= \vKoF = \sqrt{\gE\cdot\rEr} = \vKo \approx \result{\vKoP{3}{2}}
\end{align}
Das ist die erste kosmische Geschwindigkeit. In Wirklichkeit verhindert der Luftwiderstand eine solche Bahn knapp über dem Boden; Satelliten kreisen deshalb oberhalb der Atmosphäre.
